\section{Erd\H{o}s Problem \#41: independent continuation}
\noindent Date: 2026-09-23. Research attempt; no claim of a new published result.
\subsection{1) TARGET AND SOURCE AUDIT}
If $A$ is an infinite $B_3$-set (unordered triple sums, with repetitions, are unique), must $\liminf A(N)/N^{1/3}=0$?

All supplied versions considered: \texttt{41.tex}.
The source proves a valid monotone interpolation lemma and invokes an earlier no-positive-limit result. Here I prove an additional unconditional restriction on the cluster set; I do not use the inherited theorem without its proof.
\subsection{2) WORK}
Put $f(N)=A(N)N^{-1/3}$. If $t=A(N)$, the $\binom{t+2}{3}$ triple multisets have distinct integer sums between $3$ and $3N$, hence $t^3/6\le3N$ and $f(N)\le18^{1/3}$. Also $A(N+1)-A(N)\in\{0,1\}$, giving
\[
|f(N+1)-f(N)|\le (N+1)^{-1/3}
+A(N)\bigl(N^{-1/3}-(N+1)^{-1/3}\bigr)=o(1).
\]
The second term is $O(N^{-1})$ by the mean value theorem and the preceding bound.

\paragraph{The entire interval is realized.}
Write $\ell=\liminf f(N)$ and $L=\limsup f(N)$. The cluster set is exactly $[\ell,L]$. For $\ell<c<L$, there are arbitrarily late indices below $c$ and above $c$. Choose a crossing after index $M$. At its first crossing index $n$, one has $|f(n)-c|\le |f(n)-f(n-1)|$. Letting $M\to\infty$ proves that $c$ is a cluster value. The endpoints are cluster values by boundedness, and no other values can occur by the definitions of liminf and limsup.

Thus the oscillatory alternative in the source cannot have just two separated cluster values; it must realize every intervening value. This is a general small-increment observation, not a new $B_3$ density bound.
\subsection{3) ADVERSARIAL VERIFICATION}
An interval of cluster values does not supply a single convergent subsequence with successive scale ratios tending to one. Crossings of the same level may be far apart. Consequently it gives no contradiction to the supplied dense-subsequence lemma. The bound above allows repeated summands exactly as the definition requires.
\subsection{4) FINAL}
\textbf{UNRESOLVED.}
Nothing forces the left endpoint $\ell$ to vanish. The missing ingredient must use uniqueness of triple sums more strongly than the elementary counting and small-increment bounds.
\subsection{5) SOURCES CHECKED}
Read all of \texttt{41.tex}; its earlier-round Helm theorem is not needed above. Direct statement retrieval from \texttt{https://www.erdosproblems.com/41} returned 403 on 2026-09-23.
\subsection{6) COMPLETION ESTIMATE}
This is a conservative subjective measure of progress toward the original question, not a probability of correctness or a claim that the remaining work is routine.
\textbf{COMPLETION: 8\%}
